\documentclass[12pt]{article}
\usepackage[utf8]{inputenc}
\usepackage[T1]{fontenc}
\usepackage[margin=1in]{geometry}
\usepackage{setspace}
\usepackage{xcolor}
\usepackage[colorlinks=true, urlcolor=blue!60!black, linkcolor=blue!60!black]{hyperref}
\setlength{\parindent}{2em}
\setlength{\parskip}{1.25em}
\renewcommand{\baselinestretch}{1.0}
\begin{document}
\begin{center}
{\LARGE \textit{Craig v. Boren}, 429 U.S. 190 (1976)\par}
\vspace{0.8em}
{\large Case Brief --- Four-Page Edition\par}
\end{center}
\vspace{1em}

\section{Memory Jogger}

The Court struck down an Oklahoma law barring the sale of 3.2\% beer to men under 21 while permitting sales to women at 18, formally establishing ``intermediate scrutiny'' for sex-based classifications under the Equal Protection Clause (\href{https://supreme.justia.com/cases/federal/us/429/190/}{429 U.S. 190}).

\section{Detailed Case Facts}

Oklahoma law barred selling ``nonintoxicating'' 3.2\% beer to a ``minor,'' defined as a female under 18 and a male under 21. The law regulated only \textit{sale}, not possession or consumption, so an 18-year-old woman could buy beer and hand it to a male companion. In 1972, after the Tenth Circuit invalidated a sex-based criminal-responsibility age under \textit{Reed v. Reed}, Oklahoma set the age of majority at 18 for both sexes but preserved the 21/18 gap for 3.2\% beer.

Curtis Craig, a male Oklahoma State student between 18 and 21, and Carolyn Whitener, a licensed beer vendor, sued Governor Boren for declaratory and injunctive relief. Oklahoma relied on statistics linking young males to drunk driving and on the Twenty-first Amendment. Ruth Bader Ginsburg, leading the ACLU Women's Rights Project, filed an \textit{amicus} brief and sat at counsel table.

\section{Procedural History}

A three-judge district court upheld the law, finding the State's statistics showed a rational and substantial relationship to highway safety. On direct appeal, Craig turned 21, mooting his claim. The Supreme Court reversed on December 20, 1976, holding that Whitener had third-party standing.

\section{Judicial Votes}

Reversed, \textbf{7--2}.
\begin{itemize}
\item \textbf{Majority:} Justice Brennan, joined by Justices White, Marshall, Powell, and Stevens, and by Justice Blackmun except Part II-D (Twenty-first Amendment).
\item \textbf{Concurrences:} Powell; Stevens; Blackmun (in part); Stewart (in the judgment).
\item \textbf{Dissents:} Chief Justice Burger; Justice Rehnquist.
\end{itemize}

\section{Holding}

A licensed vendor directly regulated by the statute has \textit{jus tertii} standing to assert her male customers' equal protection rights. ``[C]lassifications by gender must serve important governmental objectives and must be substantially related to achievement of those objectives.'' Traffic safety is important, but arrest statistics showing 2\% of young men versus 0.18\% of young women arrested for drunk driving are too tenuous a fit, especially when the law bars only purchase. The Twenty-first Amendment does not relax equal protection scrutiny, and \textit{Goesaert v. Cleary} was disapproved insofar as it suggested otherwise.

\section{Analysis of Opinions}

\textbf{Majority (Brennan).} Whitener, forced to choose between lost sales and sanctions, was the ``obvious claimant'' for young men's rights. Synthesizing \textit{Reed}, \textit{Frontiero}, \textit{Wiesenfeld}, and \textit{Stanton}, Brennan stated the two-part test and found a 2\% correlation an ``unduly tenuous `fit.''' The surveys did not isolate 3.2\% beer, arrests might reflect police chivalry, and roadside data showed only modest gender gaps. Section~2 of the Twenty-first Amendment primarily qualifies the Commerce Clause, not individual rights.

\textbf{Powell (concurring).} Uneasy with a formal middle tier, he preferred \textit{Reed}'s ``fair and substantial relation'' test; a gap ``so easily circumvented as to be virtually meaningless'' fails it.

\textbf{Stevens (concurring).} Rejecting tiers, he found the line rested on an ``accident of birth'' and imposed a disability on all young men for a minimal deterrent effect.

\textbf{Blackmun and Stewart.} Blackmun joined all but Part II-D. Stewart concurred in the judgment, finding the law totally irrational under \textit{Reed}.

\textbf{Burger (dissenting).} Young men could sue themselves, so vendor standing was unwarranted, and only rational-basis review applied to a law that was at most ``a bit foolish.''

\textbf{Rehnquist (dissenting).} Heightened scrutiny was unjustified for laws burdening men; the new test came ``out of thin air'' with terms ``so diaphanous and elastic'' as to invite policymaking. The relevant question is relative incidence (2\% versus 0.18\%), which made the law rational.

\section{Examples: Future Applications}

\begin{enumerate}
\item \textbf{Same Side: Motorcycle Licenses.} A state sets the motorcycle license age at 21 for men and 18 for women, citing fatality data though under 1.5\% of young male riders crash. Maleness is an overbroad proxy when gender-neutral measures exist.
\item \textbf{Same Side: Nursing Scholarships.} Male scholarship recipients owe five years of service, women two, based on aggregate career-change data. Statistical generalizations fail intermediate scrutiny.
\item \textbf{Opposite Side: Uniform Drinking Age.} A state sets 21 as the drinking age for everyone. As footnote 24 noted, a gender-neutral cutoff triggers only rational-basis review.
\item \textbf{Opposite Side: Prenatal Screening.} A statute provides free prenatal screening to pregnant patients. The line tracks an actual biological difference rather than a stereotype.
\item \textbf{Fence-Sitter: Insurance Rates.} A regulator permits higher premiums for young male drivers based on actuarial data showing nearly twice the claims per mile. The correlation is stronger and financial, but \textit{Craig} warns against sex-based proxies when neutral ones exist.
\end{enumerate}

\section{Critique}

\textbf{Statistical and originalist objections.} Rehnquist exposed a real weakness: focusing on the low absolute arrest rate rather than the relative disparity proved too much, since age-based drinking laws would also fail. Scholars note Brennan's objection was really normative. Originalists argue the Equal Protection Clause contains no tiered architecture, was not understood in 1868 to bar sex classifications, and that ``intermediate scrutiny'' enacted much of the unratified Equal Rights Amendment by judicial fiat.

\textbf{Progressive and feminist objections.} Reva Siegel argues that beneath the ``faux precision'' of intermediate scrutiny lies confusion about whether sex discrimination is less grave than race discrimination. Settling on a middle tier in a low-stakes beer case institutionalized a symmetrical anti-classification rule and left women with weaker protection than strict scrutiny, a gap only partly narrowed by \textit{United States v. Virginia}.

\section{Key Quotations}

\begin{quote}
``To withstand constitutional challenge, previous cases establish that classifications by gender must serve important governmental objectives and must be substantially related to achievement of those objectives.''
\end{quote}
\noindent --- Justice Brennan, Opinion of the Court, 429 U.S. at 197.

\begin{quote}
``It is unrealistic to expect either members of the judiciary or state officials to be well versed in the rigors of experimental or statistical technique. But this merely illustrates that proving broad sociological propositions by statistics is a dubious business, and one that inevitably is in tension with the normative philosophy that underlies the Equal Protection Clause.''
\end{quote}
\noindent --- Justice Brennan, Opinion of the Court, 429 U.S. at 204.

\begin{quote}
``There is only one Equal Protection Clause. It requires every State to govern impartially. It does not direct the courts to apply one standard of review in some cases and a different standard in other cases.''
\end{quote}
\noindent --- Justice Stevens, concurring, 429 U.S. at 211--212.

\end{document}
